\chapter{Database Documentation}

\section{Database Setup}

\subsection{Option A: Local PostgreSQL}
\begin{enumerate}
    \item Install PostgreSQL locally.
    \item Create a database.
    \item Update the \texttt{.env} file with your local credentials.
\end{enumerate}

\subsection{Option B: NeonDB (Recommended)}
\begin{enumerate}
    \item Sign up at \url{https://neon.tech}.
    \item Create a new project.
    \item Copy the connection string.
    \item Update the \texttt{.env} file with NeonDB credentials.
\end{enumerate}

\section{Database Operations}

\begin{lstlisting}[language=bash, caption=Database Commands]
# Run migrations
npm run db:migrate

# Seed database
npm run db:seed

# Reset database
npm run db:reset
\end{lstlisting}

\section{Database Configuration}

\subsection{Sequelize Setup}
The application uses Sequelize ORM with PostgreSQL for database interaction. 
Key configuration is stored in \texttt{config/database.js}:

\begin{lstlisting}[language=JavaScript, caption=Sequelize Configuration]
const sequelize = new Sequelize(
  process.env.DB_NAME,
  process.env.DB_USER,
  process.env.DB_PASSWORD,
  {
    host: process.env.DB_HOST,
    port: process.env.DB_PORT,
    dialect: 'postgres',
    logging: process.env.NODE_ENV === 'development',
    dialectOptions: {
      ssl: {
        require: true,
        rejectUnauthorized: false
      }
    }
  }
);
\end{lstlisting}

\subsection{Models}
\begin{itemize}
    \item \textbf{User:} Stores authentication and profile data.
    \item Easily extensible for additional models as required.
\end{itemize}

\section{Development Workflow}

\subsection{Starting Development}
\begin{lstlisting}[language=bash]
# Install dependencies
npm install

# Set up environment
cp env.example .env
# Edit .env with your database credentials

# Start development server
npm run dev
\end{lstlisting}

\subsection{Database Changes}
\begin{lstlisting}[language=bash]
# Enable database sync for development
ENABLE_DB_SYNC=true

# Use migrations for production
npm run db:migrate
\end{lstlisting}

\subsection{Testing APIs}
\begin{itemize}
    \item Visit \url{http://localhost:3000/api-docs} for interactive documentation.
    \item Use Postman or curl for API testing.
\end{itemize}

\section{Production Deployment}

\subsection{Environment Setup}
\begin{lstlisting}[language=bash]
NODE_ENV=production
ENABLE_DB_SYNC=false  # Never sync in production
\end{lstlisting}

\subsection{Security Considerations}
\begin{itemize}
    \item Change all default secrets.
    \item Use strong JWT secrets.
    \item Enable HTTPS.
    \item Configure proper CORS origins.
    \item Set up rate limiting.
\end{itemize}

\subsection{Database}
\begin{itemize}
    \item Use production-grade PostgreSQL.
    \item Set up automated backups.
    \item Configure connection pooling.
    \item Always use migrations instead of sync.
\end{itemize}

\section{Troubleshooting}

\subsection{Common Issues}
\begin{description}
    \item[Database Connection Failed:] Verify \texttt{.env} credentials and network connectivity.
    \item[Port Already in Use:] Change the port or use \texttt{npx kill-port 3000}.
    \item[Module Not Found:] Run \texttt{npm install} and check Node.js version.
    \item[Database Sync Issues:] Set \texttt{ENABLE\_DB\_SYNC=true} for development; use migrations for production.
\end{description}

\subsection{Logs}
\begin{itemize}
    \item Application logs appear in the console.
    \item Database queries are logged in development mode.
    \item Error details are displayed in development mode.
\end{itemize}

\section{Scripts Reference}
\begin{lstlisting}[language=json, caption=Package Scripts]
{
  "start": "node server.js",
  "dev": "nodemon server.js",
  "test": "jest",
  "db:migrate": "sequelize-cli db:migrate",
  "db:seed": "sequelize-cli db:seed:all",
  "db:reset": "sequelize-cli db:drop && sequelize-cli db:create && sequelize-cli db:migrate && sequelize-cli db:seed:all"
}
\end{lstlisting}

\section{Next Steps}
\begin{itemize}
    \item Add more models as required.
    \item Implement unit and integration tests.
    \item Introduce structured logging.
    \item Add health checks and monitoring.
    \item Set up a CI/CD pipeline for automated deployments.
\end{itemize}

\section{Support}
\begin{itemize}
    \item Refer to the troubleshooting section for common issues.
    \item Review API documentation at \url{/api-docs}.
    \item Check console logs for error details.
    \item Ensure all environment variables are correctly set.
\end{itemize}


\section{Database Schema}

This section describes the schema used to store application data. Below is an example of the \texttt{notifications} table schema from our Neon database instance.

\subsection{Notifications Table}

\begin{table}[H]
\centering
\begin{tabular}{|l|l|l|l|}
\hline
\textbf{Column} & \textbf{Data Type} & \textbf{Constraints} & \textbf{Description} \\
\hline
\texttt{id} & SERIAL & PRIMARY KEY & Unique identifier for each notification. \\
\hline
\texttt{user\_id} & UUID & NOT NULL & References the user receiving the notification. \\
\hline
\texttt{message} & VARCHAR(255) & NOT NULL & The content of the notification. \\
\hline
\texttt{title} & VARCHAR(255) & NOT NULL & Title or subject of the notification. \\
\hline
\texttt{read} & BOOLEAN & NOT NULL, DEFAULT FALSE & Indicates if the notification has been read. \\
\hline
\texttt{created\_at} & TIMESTAMP WITH TIME ZONE & NOT NULL & Timestamp when the notification was created. \\
\hline
\end{tabular}
\caption{Schema for the \texttt{notifications} table.}
\end{table}




\section{Choice Justification}

In selecting a database solution for this project, Neon DB emerged as the most effective and forward-thinking option, surpassing traditional relational databases in flexibility, scalability, and performance. Neon offers a fully managed, serverless PostgreSQL platform, meaning it provides all the robust relational features of PostgreSQL while introducing modern cloud-native benefits such as autoscaling, high availability, and instant branching for development and testing. 

Our project required a solution that could elegantly handle relational data, linking entities such as users, resources, and reports without compromising speed or reliability. Neon was the ideal fit because it allowed us to seamlessly model relationships between tables while benefiting from PostgreSQL’s mature query capabilities. Unlike many other databases that either limit relational functionality (as seen in certain NoSQL systems) or impose heavy maintenance overhead (as with self-hosted relational databases), Neon combines ease of use with enterprise-grade reliability. 

Additionally, Neon’s separation of compute and storage makes scaling cost-efficient and straightforward, ensuring that our application can evolve without disruptive migrations or manual scaling efforts. Its seamless integration with ORMs like Sequelize further streamlines development, allowing our team to maintain a clean codebase while interacting with the database intuitively. This strategic choice ensures that we can confidently build a future-proof application with high performance, robust data integrity, and operational simplicity.


\section{High End Overview}

\begin{itemize}
    \item \textbf{User-Centric Relationships}: 
    \begin{itemize}
        \item The \texttt{User} table is central; most tables reference it via \texttt{user\_id}, establishing clear ownership.
        \item Tables like \texttt{Notifications}, \texttt{PrivateChats}, \texttt{Progress}, and \texttt{Study\_sessions} are one-to-many with \texttt{User}.
    \end{itemize}

    \item \textbf{Many-to-Many Relationships}:
    \begin{itemize}
        \item \texttt{UserCourses} links \texttt{Users} and \texttt{Courses}.
        \item \texttt{Group\_members} links \texttt{Users} and \texttt{Study\_groups}.
        \item \texttt{Likes} links \texttt{Users} and \texttt{Resources}.
        \item These join tables avoid data duplication but can grow large with many users or resources.
    \end{itemize}

    \item \textbf{Self-Referencing Relationships}:
    \begin{itemize}
        \item \texttt{Follows} and \texttt{Follows\_requests} implement social connections between users.
        \item Requires careful indexing to maintain query performance as the user count grows.
    \end{itemize}

    \item \textbf{Resource-Centric Relationships}:
    \begin{itemize}
        \item \texttt{Resources} connect to \texttt{Resource\_tags}, \texttt{Resource\_threads}, and \texttt{Resource\_reports}.
        \item High-traffic resources may create join hotspots.
    \end{itemize}
    
     \item \textbf{Structure observation}:
    \begin{itemize}
        \item \textbf{Normalisation}--- our data is clearly normalised as a result of non redundant data each piece of information is stored only once (e.g., UserCourses links Users and Courses instead of repeating course info for every user).Separate tables for entities-  Users, Courses, Resources, Likes, Groups, etc., each have their own table.
        \item Foreign keys instead of repeated data – Relationships are handled via foreign keys (user\_id, course\_id, etc.) rather than duplicating user or course details.
    \end{itemize}

    \item \textbf{Scalability Considerations}:
    \begin{itemize}
        \item Proper indexing on foreign keys is essential (\texttt{user\_id}, \texttt{resource\_id}, \texttt{course\_id}, etc.).
        \item Many-to-many relationships scale quadratically, with active users and resources.
        \item Database scalability is addressed with indexing, query optimization, caching (e.g., Redis).
    \end{itemize}
\end{itemize}
